\documentclass{article}
\usepackage{../assignment_preamble}

\title{Homework 1}
\author{Ravi Kini}
\date{October 6, 2025}

\begin{document}

\maketitle

\problem{}
\subproblem{(a)}
\begin{align*}
    \mathrm{col} \begin{bmatrix}
        1 & 0 & 0 \\
        0 & 1 & 0 \\
        0 & 0 & 0 \\
    \end{bmatrix}
\end{align*}
The matrix is already in row echelon form. Since there are $2$ pivot columns, the column space has dimension $2$.

\subproblem{(b)}
\begin{align*}
    \mathrm{span}\{(1, 1, 1), (1, -1, 1), (-1, 1, 1), (1, 1, -1)\}
\end{align*}
The dimension of the span is equal to the dimension of:
\begin{align*}
    \mathrm{col} \begin{bmatrix}
        1 & 1 & -1 & 1 \\
        1 & -1 & 1 & 1 \\
        1 & 1 & 1 & -1 \\
    \end{bmatrix}
\end{align*}
Putting the matrix in row echelon form using Gaussian elimination:
\begin{align*}
    & \begin{gmatrix}[b]
        1 & 1 & -1 & 1 \\
        1 & -1 & 1 & 1 \\
        1 & 1 & 1 & -1
    \rowops
    \add[-1][]{0}{1}
    \add[-1][]{0}{2}
    \end{gmatrix}  \\
    & \begin{gmatrix}[b]
        1 & 1 & -1 & 1 \\
        0 & -2 & 2 & 0 \\
        0 & 0 & 2 & -2
    \end{gmatrix} 
\end{align*}
Since there are $3$ pivot columns, the column space and the span have dimension $3$.

\subproblem{(c)}
\begin{align*}
    \mathrm{span}\{(2, 7, 9), (3, 5, 1), (0, 1, 0)\}
\end{align*}
The dimension of the span is equal to the dimension of:
\begin{align*}
    \mathrm{col} \begin{bmatrix}
        2 & 3 & 0 \\
        7 & 5 & 1 \\
        9 & 1 & 0 \\
    \end{bmatrix}
\end{align*}
Putting the matrix in row echelon form using Gaussian elimination:
\begin{align*}
    & \begin{gmatrix}[b]
        2 & 3 & 0 \\
        7 & 5 & 1 \\
        9 & 1 & 0
    \rowops
    \add[-\frac{7}{2}][]{0}{1}
    \add[-\frac{9}{2}][]{0}{2}
    \end{gmatrix}  \\
    & \begin{gmatrix}[b]
        2 & 3 & 0 \\
        0 & -\frac{11}{2} & 1 \\
        0 & -\frac{19}{2} & 0
    \rowops
    \add[\frac{19}{11}][]{1}{2}
    \end{gmatrix}   \\
    & \begin{gmatrix}[b]
        2 & 3 & 0 \\
        0 & -\frac{11}{2} & 1 \\
        0 & 0 & \frac{19}{11}
    \end{gmatrix} 
\end{align*}
Since there are $3$ pivot columns, the column space and the span have dimension $3$.

\subproblem{(d)}
\begin{align*}
    \mathrm{col} \begin{bmatrix}
        1 & 1 & 0 \\
        1 & 1 & 0 \\
        0 & 0 & 1 \\
    \end{bmatrix}
\end{align*}
Putting the matrix in row echelon form using Gaussian elimination:
\begin{align*}
    & \begin{gmatrix}[b]
        1 & 1 & 0 \\
        1 & 1 & 0 \\
        0 & 0 & 1
    \rowops
    \add[-1][]{0}{1}
    \end{gmatrix}  \\
    & \begin{gmatrix}[b]
        1 & 1 & 0 \\
        0 & 0 & 0 \\
        0 & 0 & 1
    \rowops
    \swap{1}{2}
    \end{gmatrix}   \\
    & \begin{gmatrix}[b]
        1 & 1 & 0 \\
        0 & 0 & 1 \\
        0 & 0 & 0
    \end{gmatrix} 
\end{align*}
Since there are $2$ pivot columns, the column space has dimension $2$.
\clearpage

\problem{}
\subproblem{(a)}
\begin{align*}
    f(x, y, z) = 0
\end{align*}
A linear function must preserve vector addition.
\begin{align*}
    f(x_1, y_1, z_1) + f(x_2, y_2, z_2) & = 0 + 0 = 0 \\
    f(x_1 + x_2, y_1 + y_2, z_1 + z_2) & = 0 \\
\end{align*}
A linear function must preserve scalar multiplication. For $a \in \mathbb{R}$:
\begin{align*}
    a \cdot f(x, y, z) & = a \cdot 0 = 0 \\
    f(ax, ay, az) & = 0
\end{align*}
Since both vector addition and scalar multiplication are preserved, the function is linear.

\subproblem{(b)}
\begin{align*}
    f(x, y, z) = 1
\end{align*}
A linear function must preserve vector addition.
\begin{align*}
    f(x_1, y_1, z_1) + f(x_2, y_2, z_2) & = 1 + 1 = 2 \\
    f(x_1 + x_2, y_1 + y_2, z_1 + z_2) & = 1 \\
\end{align*}
A linear function must preserve scalar multiplication. For $a \in \mathbb{R}$:
\begin{align*}
    a \cdot f(x, y, z) & = a \cdot 1 = a \\
    f(ax, ay, az) & = 1
\end{align*}
Since neither vector addition nor scalar multiplication are preserved, the function is not linear.

\subproblem{(c)}
\begin{align*}
    f(x, y, z) = (1 + x, 2z)
\end{align*}
A linear function must preserve vector addition.
\begin{align*}
    f(x_1, y_1, z_1) + f(x_2, y_2, z_2) & = (1 + x_1, 2z_1) + (1 + x_2, 2z_2) \\
    & = (2 + x_1 + x_2, 2z_1 + 2z_2) \\
    f(x_1 + x_2, y_1 + y_2, z_1 + z_2) & = (1 + (x_1 + x_2), 2(z_1 + z_2)) \\
    & = (1 + x_1 + x_2, 2z_1 + 2z_2)
\end{align*}
A linear function must preserve scalar multiplication. For $a \in \mathbb{R}$:
\begin{align*}
    a \cdot f(x, y, z) & = a \cdot (1 + x, 2z) = (a(1 + x), 2az) \\
    & = (a + ax, 2az) \\
    f(ax, ay, az) & = (1 + ax, 2az)
\end{align*}
Since neither vector addition nor scalar multiplication are preserved, the function is not linear.

\subproblem{(d)}
\begin{align*}
    f(x, y, z) = (x, 2x)
\end{align*}
A linear function must preserve vector addition.
\begin{align*}
    f(x_1, y_1, z_1) + f(x_2, y_2, z_2) & = (x_1, 2x_1) + (x_2, 2x_2) \\
    & = (x_1 + x_2, 2x_1 + 2x_2) \\
    f(x_1 + x_2, y_1 + y_2, z_1 + z_2) & = (x_1 + x_2, 2(x_1 + x_2)) \\
    & = (x_1 + x_2, 2x_1 + 2x_2)
\end{align*}
A linear function must preserve scalar multiplication. For $a \in \mathbb{R}$:
\begin{align*}
    a \cdot f(x, y, z) & = a \cdot (x, 2x) = (ax, 2ax) \\
    f(ax, ay, az) & = (ax, 2ax)
\end{align*}
Since both vector addition and scalar multiplication are preserved, the function is linear.


\subproblem{(e)}
\begin{align*}
    f(x, y, z) = (2x + 3y, x, 0)
\end{align*}
A linear function must preserve vector addition.
\begin{align*}
    f(x_1, y_1, z_1) + f(x_2, y_2, z_2) & = (2x_1 + 3y_1, x_1, 0) + (2x_2 + 3y_2, x_2, 0) \\
    & = (2x_1 + 2x_2 + 3y_1 + 3y_2, x_1 + x_2, 0) \\
    f(x_1 + x_2, y_1 + y_2, z_1 + z_2) & = (2(x_1 + x_2) + 3(y_1 + y_2), x_1 + x_2, 0) \\
    & = (2x_1 + 2x_2 + 3y_1 + 3y_2, x_1 + x_2, 0)
\end{align*}
A linear function must preserve scalar multiplication. For $a \in \mathbb{R}$:
\begin{align*}
    a \cdot f(x, y, z) & = a \cdot (2x + 3y, x, 0) = (a(2x + 3y), x, 0) \\
    & = (2ax + 3ay, ax, 0) \\
    f(ax, ay, az) & = (2ax + 3ay, ax, 0)
\end{align*}
Since both vector addition and scalar multiplication are preserved, the function is linear.
\clearpage

\problem{}
Let $\mathcal{U}_1, \mathcal{U}_2$ be subspaces of vector space $\mathcal{V}$. Let $u_1, u_2 \in \mathcal{U}_1 \cap \mathcal{U}_2$. Since $u_1, u_2 \in \mathcal{U}_1$, $u_1 + u_2 \in \mathcal{U}_1$. Similarly, $u_1 + u_2 \in \mathcal{U}_2$. Consequently, $u_1 + u_2 \in \mathcal{U}_1 \cap \mathcal{U}_2$, and vector addition is preserved. Let $a \in \mathbb{R}$. Since $u_1 \in \mathcal{U}_1$, $a \cdot u_1 \in \mathcal{U}_1$. Similarly, $a \cdot u_1 \in \mathcal{U}_2$. Consequently, $a \cdot u_1 \in \mathcal{U}_1 \cap \mathcal{U}_2$, and scalar multiplication is preserved. Since $\mathcal{U}_1 \cap \mathcal{U}_2$ is closed under both vector addition and scalar multiplication, $\mathcal{U}_1 \cap \mathcal{U}_2$ is a subspace of $\mathcal{V}$.

However, ${U}_1 \cup \mathcal{U}_2$ is not necessarily a subspace of $\mathcal{V}$. Let $\mathcal{U}_1 = \mathrm{span}\{(1, 0)\}$ and $\mathcal{U}_2 = \mathrm{span}\{(0, 1)\}$, which are subspaces of $\mathcal{V} = \mathbb{R}^2$. Consider $u_1 = (1, 0) \in \mathcal{U}_1$ and $u_2 = (0, 1) \in \mathcal{U}_2$. $u_1 + u_2 = (1, 1)$, but since $u_1 + u_2 \not\in \mathcal{U}_1$ and $u_1 + u_2 \not\in \mathcal{U}_2$, $u_1 + u_2 \not\in \mathcal{U}_1 \cup \mathcal{U}_2$. Since $\mathcal{U}_1 \cup \mathcal{U}_2$ is not closed under vector addition, $\mathcal{U}_1 \cup \mathcal{U}_2$ is not a subspace of $\mathcal{V}$.
\clearpage

\problem{}
\begin{align*}
    A = \begin{bmatrix}
        r_1^T \\
        r_2^T \\
        \vdots \\
        r_m^T 
    \end{bmatrix} = \begin{bmatrix}
        c_1 & c_2 & \cdots & c_n 
    \end{bmatrix}
\end{align*}
The elements of $A^TA$ can be expressed as:
\begin{align*}
    (A^TA)_{ij} & = (A^T)_{i.} \cdot (A)_{.j} \\
    & = c_i \cdot c_j
\end{align*}
The elements of $AA^T$ can be expressed as:
\begin{align*}
    (AA^T)_{ij} & = (A)_{i.} \cdot (A^T)_{.j} \\
    & = r_i \cdot r_j
\end{align*}
\clearpage

\problem{}
Let $A, B \in \mathbb{R}^{m\times n}$ be invertible. Then:
\begin{align*}
    A^T(A^{-1})^T & = (A^{-1}A)^T \\
    & = I^T = I \\
    & = A^T(A^T)^{-1} \\
    (A^{-1})^T & = (A^T)^{-1}
\end{align*}
Further:
\begin{align*}
    (AB)(B^{-1}A^{-1}) & = A(BB^{-1})A^{-1} \\
    & = AIA^{-1} = AA^{-1} \\
    & = I \\
    & = (AB)(AB)^{-1} \\
    (B^{-1}A^{-1}) & = (AB)^{-1}
\end{align*}
\clearpage

\problem{}
\subproblem{(a)}
Let $B \in \mathbb{R}^{m\times k}$ such that $B^TB$ is invertible. Then:
\begin{align*}
    (B(B^TB)^{-1}B^T)^2 & = (B(B^TB)^{-1}B^T)(B(B^TB)^{-1}B^T) \\
    & = B(B^TB)^{-1}(B^TB)(B^TB)^{-1}B^T \\
    & = B(B^TB)^{-1}IB^T = B(B^TB)^{-1}B^T \\
\end{align*}
Since $(B(B^TB)^{-1}B^T)^2 = B(B^TB)^{-1}B^T$, $B(B^TB)^{-1}B^T$ is idempotent.

\subproblem{(b)}
Let $A$ be idempotent. Then:
\begin{align*}
    (I - A)^2 & = I^2 - IA - AI + A^2 \\
    & = I - 2A + A = I - A
\end{align*}
Since $(I - A)^2 = I - A$, $I - A$ is idempotent.

\subproblem{(c)}
Let $A$ be idempotent. Consider $\frac{1}{2}I - A$ and $2I - 4A$. Then:
\begin{align*}
    \left(\frac{1}{2}I - A\right)(2I - 4A) & = I^2 - 2AI - 2IA + 4A^2 \\
    & = I - 4A + 4A = I
\end{align*}
Since $\left(\frac{1}{2}I - A\right)(2I - 4A) = I$, $\frac{1}{2}I - A$ is invertible, with an inverse of $2I - 4A$.

\subproblem{(d)}
Let $A$ be idempotent. Let $x \neq 0$ and $\lambda \in \mathbb{R}$ such that $Ax = \lambda x$. Then:
\begin{align*}
    Ax & = \lambda x \\
    A^2x = Ax & = \lambda Ax \\
    \lambda x & = \lambda(\lambda x) = \lambda^2 x \\
    (\lambda^2 - \lambda)x & = \lambda(\lambda - 1)x = 0 \\
    \lambda & = 0, 1
\end{align*}
\clearpage

\problem{}
\begin{align*}
    \sin(x) = x - \frac{x^3}{3!} + \frac{x^5}{5!} - \frac{x^7}{7!} + \ldots
\end{align*}
\subproblem{(a)}
\begin{align*}
    \sin(x) \approx x
\end{align*}
The relative forward error at $x$ is $\left|\frac{\sin(x) - x}{\sin(x)}\right|$. The relative backward error is $\left|\frac{\frac{x^3}{6}}{x}\right|$.

\begin{center}
\begin{tabular}{|c|c|c|}
    \hline
    $x$ & rel. forward err. & rel. backward err. \\
    \hline
    $0.1$ & $0.00167$ & $0.00167$ \\
    $0.5$ & $0.04291$ & $0.04167$ \\
    $1.0$ & $0.18840$ & $0.16667$ \\
    \hline
\end{tabular}
\end{center}

\subproblem{(b)}
\begin{align*}
    \sin(x) \approx x - \frac{x^3}{6}
\end{align*}
The relative forward error at $x$ is $\left|\frac{\sin(x) - (x - \frac{x^3}{6})}{\sin(x)}\right|$. The relative backward error is $\left|\frac{\frac{x^3}{120}}{x}\right|$.

\begin{center}
\begin{tabular}{|c|c|c|}
    \hline
    $x$ & rel. forward err. & rel. backward err. \\
    \hline
    $0.1$ & $8.3452 \times 10^{-7}$ & $8.3333 \times 10^{-7}$ \\
    $0.5$ & $0.00054$ & $0.00052$ \\
    $1.0$ & $0.00967$ & $0.00833$ \\
    \hline
\end{tabular}
\end{center}
\clearpage

\problem{}
\begin{align*}
    f(x) = \log x
\end{align*}
\subproblem{(a)}
The condition number $\kappa_f(x)$ is:
\begin{align*}
    \kappa_f(x) & = \left|\frac{x \cdot f'(x)}{x}\right| \\
    & = \left|\frac{x \cdot \frac{1}{x}}{\log x}\right| = \left|\frac{1}{\log x}\right|
\end{align*}

\subproblem{(b)}
The relative forward error at $x$ is $\left|\frac{f(x + \Delta x) - f(x)}{f(x)}\right|$. The relative backward error is $\left|\frac{\Delta x}{x}\right|$. Using $\Delta x = 10^{-4}$:

\begin{center}
\begin{tabular}{|c|c|c|c|}
    \hline
    $x$ & $f(x)$ & $f(x + \Delta x)$ & $\kappa_f(x)$ \\
    \hline
    $1.005$ & $0.00498$ & $0.00508$ & $200.49958$ \\
    $1.004$ & $0.00399$ & $0.00409$ & $250.49967$ \\
    $1.003$ & $0.00300$ & $0.00310$ & $333.83308$ \\
    $1.002$ & $0.00200$ & $0.00210$ & $500.49983$ \\
    $1.001$ & $0.00100$ & $0.00110$ & $1000.49992$ \\
    \hline
\end{tabular}
\end{center}
\begin{center}
\begin{tabular}{|c|c|c|c|}
    \hline
    $x$ & rel. forward err. & rel. backward err. & ratio  \\
    \hline
    $1.005$ & $0.01994$ & $0.00010$ & $200.48961$  \\
    $1.004$ & $0.02495$ & $0.00010$ & $250.48719$ \\
    $1.003$ & $0.03328$ & $0.00010$ & $333.81644$ \\
    $1.002$ & $0.049948$ & $0.00010$ & $500.47486$\\
    $1.001$ & $0.09995$ & $0.00010$ & $1000.44995$ \\
    \hline
\end{tabular}
\end{center}
The relative backward error stays fairly constant (close to $10^{-4}$). However, since $|\log x|$ grows dramatically as $x$ approaches $1$, the condition number grows dramatically as well, resulting in a large upper bound on the relative forward error and a much larger relative change in $f(x + \Delta x)$ in comparison to $f(x)$.
\end{document}